\documentclass{article}

\begin{document}

\section{Introduction}
Like the sectional units addressed in Section \ref{sec:section_units}, LaTeX assigns serial numbers to many environments or elements of an environment (e.g., table, figure, equation, or \texttt{\textbackslash item} as discussed in the following sections). This default numbering system eliminates the possibility of committing any mistake as may happen in manual numbering.

Moreover, \LaTeX\ allows you to label a numbered item by a unique reference key, which can be used to refer to the item in any part within the same document (unnumbered items, like \texttt{\textbackslash paragraph\{\}}, cannot be referred to in this way). The labeling and referring of an item are performed through \texttt{\textbackslash label\{rkey\}} and \texttt{\textbackslash ref\{rkey\}} respectively, where \texttt{rkey} is the assigned unique reference key of the item\textsuperscript{2}.

\subsection{Sectional Units}
\label{sec:section_units}
In this section, we discuss the use of sectional units in LaTeX.

\end{document}
